\documentclass[11pt,a4paper]{article}

% ── Packages ──────────────────────────────────────────────────────────────
\usepackage[utf8]{inputenc}
\usepackage[T1]{fontenc}
\usepackage{amsmath,amssymb,amsthm}
\usepackage{mathtools}
\usepackage{bm}
\usepackage{geometry}
\geometry{margin=1in}
\usepackage{hyperref}
\hypersetup{colorlinks=true,linkcolor=blue,citecolor=blue,urlcolor=blue}
\usepackage{listings}
\usepackage{xcolor}
\usepackage{titlesec}
\usepackage{enumitem}
\usepackage{booktabs}
\usepackage{fancyhdr}

% ── Code listing style ───────────────────────────────────────────────────
\definecolor{codebg}{RGB}{248,248,248}
\definecolor{codeframe}{RGB}{200,200,200}
\definecolor{codecomment}{RGB}{0,128,0}
\definecolor{codestring}{RGB}{163,21,21}
\definecolor{codekeyword}{RGB}{0,0,200}

\lstdefinestyle{pythonstyle}{
    language=Python,
    backgroundcolor=\color{codebg},
    frame=single,
    rulecolor=\color{codeframe},
    basicstyle=\ttfamily\small,
    keywordstyle=\color{codekeyword}\bfseries,
    commentstyle=\color{codecomment}\itshape,
    stringstyle=\color{codestring},
    showstringspaces=false,
    breaklines=true,
    breakatwhitespace=false,
    tabsize=4,
    captionpos=b,
    numbers=left,
    numberstyle=\tiny\color{gray},
    numbersep=8pt,
    xleftmargin=1.5em,
    framexleftmargin=1.5em,
    aboveskip=1em,
    belowskip=1em,
    morekeywords={self,True,False,None,as,with,lambda},
    literate={→}{$\rightarrow$}1 {←}{$\leftarrow$}1,
}
\lstset{style=pythonstyle}

% ── Header / Footer ──────────────────────────────────────────────────────
\pagestyle{fancy}
\fancyhf{}
\rhead{Burgers' Equation PINN --- PyTorch Implementation}
\lhead{Optimizing the Optimizer}
\cfoot{\thepage}

% ── Title ─────────────────────────────────────────────────────────────────
\title{%
  \textbf{Line-by-Line Analysis of the Burgers' Equation PINN}\\[4pt]
  \large PyTorch Implementation with Self-Scaled Broyden Optimization\\[4pt]
  \normalsize Based on \emph{Optimizing the Optimizer for PINNs} (Kiyani et al., CMAME 2025)
}
\author{}
\date{\today}

% ══════════════════════════════════════════════════════════════════════════
\begin{document}
\maketitle
\tableofcontents
\newpage

% ======================================================================
\section{Overview}
% ======================================================================

This document provides a complete, section-by-section explanation of the
\texttt{Burgers\_Equation\_Pytorch.ipynb} notebook. The notebook solves
the one-dimensional viscous Burgers' equation with a \textbf{Physics-Informed
Neural Network (PINN)}, using a two-phase training strategy:

\begin{enumerate}
  \item \textbf{Phase 1 --- Adam} (first-order): 1\,000 warmup iterations to
        bring the network parameters into a good basin.
  \item \textbf{Phase 2 --- Self-Scaled Broyden} (second-order, quasi-Newton):
        10\,000 iterations using a patched \texttt{scipy.optimize.minimize}
        that implements the self-scaled Broyden family updates from the paper.
\end{enumerate}

The PDE, its boundary/initial conditions, the neural network architecture,
the loss function, and both optimization phases are covered below.

% ======================================================================
\section{Section 1: Hyperparameters}
\label{sec:hyperparams}
% ======================================================================

\subsection*{The PDE}

The \textbf{1D viscous Burgers' equation} is:
\begin{equation}
\boxed{
  \frac{\partial u}{\partial t} + u\,\frac{\partial u}{\partial x}
  = \nu\,\frac{\partial^2 u}{\partial x^2}
}
\label{eq:burgers}
\end{equation}
where $u(t,x)$ is the unknown velocity field and $\nu$ is the kinematic
viscosity.  The nonlinear convection term $u\,u_x$ causes wave steepening,
while the diffusion term $\nu\,u_{xx}$ smooths gradients.  When $\nu$ is
small (as here, $\nu = 0.01/\pi \approx 0.00318$), a steep \emph{shock-like}
front develops, making the problem challenging for neural networks.

\subsection*{Domain and conditions}

\begin{align}
  &\text{Spatial domain:}  &x &\in [-1,\, 1] \\
  &\text{Time domain:}     &t &\in [0,\, 1]   \\
  &\text{Initial condition:} &u(0, x) &= -\sin(\pi x) \label{eq:ic}\\
  &\text{Boundary conditions:} &u(t, -1) = u(t, 1) &= 0 \quad\text{(Dirichlet)} \label{eq:bc}
\end{align}

\subsection*{Network architecture}

The MLP has the layer structure:
\[
  (t,x) \;\xrightarrow{\;\text{Linear}\;}\;
  \underbrace{20 \to 20 \to 20 \to 20}_{\text{4 hidden layers, tanh}}
  \;\xrightarrow{\;\text{Linear}\;}\; u
\]
giving \texttt{layer\_dims = (2, 20, 20, 20, 20, 1)}.

\subsection*{Training schedule}

\begin{center}
\begin{tabular}{lrl}
\toprule
\textbf{Parameter} & \textbf{Value} & \textbf{Meaning} \\
\midrule
\texttt{Nepochs\_ADAM} & 1\,000   & Adam warmup iterations \\
\texttt{Nchange}       & 500      & Resample collocation points every 500 iters \\
\texttt{Nint}          & 8\,000   & Interior collocation points \\
\texttt{N0}            & 500      & Initial-condition points ($t{=}0$) \\
\texttt{Nb}            & 500      & Boundary points ($x{=}{\pm}1$) \\
\texttt{$k_1, k_2$}   & 1, 1     & RAD parameters \\
$\nu$                  & $0.01/\pi$ & Viscosity \\
\bottomrule
\end{tabular}
\end{center}

\subsection*{Code}

\begin{lstlisting}
scaling = 1
L1 = 20
layer_dims = (2, L1, L1, L1, L1, 1)

L = 2
x0 = -1
xf = x0 + L   # = 1
tfinal = 1
nu = 0.01 / np.pi

Nepochs_ADAM = 1000
Nchange = 500
Nint    = 8000
N0      = 500
Nb      = 500
Nprint  = 100

k1 = 1;  k2 = 1
rad_args = (k1, k2)
\end{lstlisting}


% ======================================================================
\section{Section 2: Neural Network Definition}
\label{sec:network}
% ======================================================================

\subsection*{Mathematical description}

A fully-connected feed-forward neural network (MLP) maps $(t,x)$ to $\hat{u}$:
\begin{equation}
  \hat{u}(t,x;\bm{\theta}) = \sigma_{\text{out}}\!\Bigl(
    \bm{W}^{[5]}\,\sigma\!\bigl(
      \bm{W}^{[4]}\,\sigma\!\bigl(
        \bm{W}^{[3]}\,\sigma\!\bigl(
          \bm{W}^{[2]}\,\sigma\!\bigl(
            \bm{W}^{[1]} \begin{pmatrix} t \\ x \end{pmatrix}
            + \bm{b}^{[1]}
          \bigr) + \bm{b}^{[2]}
        \bigr) + \bm{b}^{[3]}
      \bigr) + \bm{b}^{[4]}
    \bigr) + \bm{b}^{[5]}
  \Bigr)
\label{eq:mlp}
\end{equation}
where:
\begin{itemize}[nosep]
  \item $\sigma(z) = \tanh(z)$ is the hidden-layer activation (smooth and $C^\infty$,
        essential because we differentiate through the network to get $u_t$, $u_x$, $u_{xx}$).
  \item $\sigma_{\text{out}}(z) = z \cdot \texttt{scaling}$ is the output activation
        (identity when \texttt{scaling}$=1$).
  \item $\bm{\theta} = \{(\bm{W}^{[\ell]}, \bm{b}^{[\ell]})\}_{\ell=1}^{5}$
        collects all trainable parameters.
\end{itemize}

The total parameter count is:
\[
  (2 \times 20 + 20) + 3 \times (20 \times 20 + 20) + (20 \times 1 + 1)
  = 60 + 1260 + 21 = 1341
\]

\subsection*{Weight initialization}

The final linear layer uses a \emph{variance-scaling} initializer to keep
the output variance $\approx 1$:
\begin{equation}
  n = \frac{\text{fan\_in} + \text{fan\_out}}{2}, \qquad
  \ell = \sqrt{\frac{3 \cdot s}{n}}, \qquad
  W_{ij} \sim \mathcal{U}(-\ell,\, \ell), \qquad
  b_j = 0
\label{eq:init}
\end{equation}
where $s = 1/\texttt{scaling}^2$.  This matches TensorFlow's
\texttt{VarianceScaling(mode='fan\_avg')}.

\subsection*{Code}

\begin{lstlisting}
class CustomActivation(nn.Module):
    def __init__(self, **kwargs):
        super().__init__()
    def forward(self, x):
        return x * scaling


def init_variance_scaling_fan_avg_uniform(linear, scale):
    fan_in, fan_out = linear.in_features, linear.out_features
    n = 0.5 * (fan_in + fan_out)
    limit = math.sqrt(3.0 * scale / n)
    with torch.no_grad():
        linear.weight.uniform_(-limit, limit)
        if linear.bias is not None:
            linear.bias.zero_()


class Net(nn.Module):
    def __init__(self):
        super().__init__()
        self.net = nn.Sequential(
            nn.Linear(layer_dims[0], layer_dims[1]), nn.Tanh(),
            nn.Linear(layer_dims[1], layer_dims[2]), nn.Tanh(),
            nn.Linear(layer_dims[2], layer_dims[3]), nn.Tanh(),
            nn.Linear(layer_dims[3], layer_dims[4]), nn.Tanh(),
            nn.Linear(layer_dims[4], layer_dims[5]),
            CustomActivation(),
        )
        final_linear = self.net[-2]
        init_variance_scaling_fan_avg_uniform(
            final_linear, scale=1.0/(scaling**2))

    def forward(self, x):
        return self.net(x)

N = Net()
\end{lstlisting}


% ======================================================================
\section{Section 3: Data Sampling Functions}
\label{sec:sampling}
% ======================================================================

\subsection*{Mathematical description}

The PINN loss requires three sets of collocation points sampled randomly
at each training stage:

\begin{enumerate}[nosep]
  \item \textbf{Interior points} $\{(t_i, x_i)\}_{i=1}^{N_{\mathrm{int}}}$,
        $(t_i, x_i) \sim \mathcal{U}([0,1] \times [-1,1])$, where the PDE
        residual is enforced.
  \item \textbf{Initial-condition points} $\{(0, x_j)\}_{j=1}^{N_0}$,
        $x_j \sim \mathcal{U}([-1,1])$, where the IC is enforced.
  \item \textbf{Boundary points} $\{(t_k, x_k)\}_{k=1}^{N_b}$,
        $t_k \sim \mathcal{U}([0,1])$, $x_k \in \{-1, 1\}$ chosen uniformly
        at random, where the BCs are enforced.
\end{enumerate}

Additionally, \textbf{Residual-based Adaptive Distribution (RAD)} resamples
interior points proportionally to the PDE residual magnitude, concentrating
points where the network is performing worst:
\begin{equation}
  p(t_i, x_i) \;\propto\;
  \frac{|f(t_i, x_i)|^{k_1}}{\overline{|f|^{k_1}}} + k_2
\label{eq:rad}
\end{equation}
where $f$ is the PDE residual, $\overline{(\cdot)}$ denotes the mean,
$k_1$ controls the aggressiveness of focusing, and $k_2$ ensures a
baseline probability everywhere.

\subsection*{Code}

\begin{lstlisting}
def generate_inputs(Nint):
    t = tfinal * np.random.rand(Nint)
    x = L * np.random.rand(Nint) + x0
    X = np.hstack((t[:, None], x[:, None]))
    return torch.as_tensor(X)

def initial_points(N0):
    t = np.zeros(N0)
    x = L * np.random.rand(N0) + x0
    X = np.hstack((t[:, None], x[:, None]))
    return torch.as_tensor(X)

def boundary_points(Nb):
    t = tfinal * np.random.rand(Nb)
    x = L * np.random.randint(0, 2, size=(Nb,)) + x0
    X = np.hstack((t[:, None], x[:, None]))
    return torch.as_tensor(X)

def generate_validation(Nt, Nx):
    x = np.linspace(x0, xf, Nx)
    t = np.linspace(0, tfinal, Nt)
    x, t = np.meshgrid(x, t)
    X = np.hstack((t.flatten()[:, None],
                   x.flatten()[:, None]))
    return torch.as_tensor(X), t, x

def adaptive_rad(N, Nint, rad_args, Nsampling=50000):
    Xtest = generate_inputs(Nsampling)
    k1, k2 = rad_args
    Y = torch.abs(get_results(N, Xtest)[-1]).reshape(-1)
    w = (Y ** k1)
    err_eq = w / w.mean() + k2
    p = (err_eq / err_eq.sum()).clamp_min(1e-12)
    X_ids = torch.multinomial(
        p, num_samples=Nint, replacement=False)
    return Xtest[X_ids]
\end{lstlisting}


% ======================================================================
\section{Section 4: PDE Conditions (IC and BC)}
\label{sec:conditions}
% ======================================================================

\subsection*{Mathematical description}

The exact conditions that the PINN must satisfy:

\begin{align}
  \text{Initial condition:}\quad & u(0, x) = -\sin(\pi x)
    && \forall\, x \in [-1,1] \label{eq:ic2}\\
  \text{Left BC:}\quad           & u(t, -1) = 0
    && \forall\, t \in [0,1] \label{eq:bcl}\\
  \text{Right BC:}\quad          & u(t, +1) = 0
    && \forall\, t \in [0,1] \label{eq:bcr}
\end{align}

At $t=0$, $u(0,x) = -\sin(\pi x)$ is a smooth sinusoidal profile.
As time evolves, the nonlinear convection steepens this profile
into a \emph{shock-like} front near $x=0$ for small $\nu$.

\subsection*{Code}

\begin{lstlisting}
def uinit(X):
    x = X[:, 1:2]                     # shape (N, 1)
    return -torch.sin(torch.pi * x)   # -sin(pi*x)

def uleft(X):
    Xleft = X[X[:, 1] == x0]
    t = Xleft[:, 0]
    return torch.zeros((t.shape[0], 1),
                       dtype=torch.float64)

def uright(X):
    Xright = X[X[:, 1] == xf]
    t = Xright[:, 0]
    return torch.zeros((t.shape[0], 1),
                       dtype=torch.float64)
\end{lstlisting}


% ======================================================================
\section{Section 5: Forward Pass and PDE Residual}
\label{sec:residual}
% ======================================================================

\subsection*{Mathematical description}

This is the \textbf{core of the PINN}.  Given the network's prediction
$\hat{u}(t,x;\bm\theta)$, we compute the PDE residual using
\emph{automatic differentiation}:

\begin{equation}
  f(t,x;\bm\theta) = \frac{\partial \hat{u}}{\partial t}
  + \hat{u}\,\frac{\partial \hat{u}}{\partial x}
  - \nu\,\frac{\partial^2 \hat{u}}{\partial x^2}
\label{eq:residual}
\end{equation}

The derivatives are computed by \texttt{torch.autograd.grad}:

\begin{enumerate}[nosep]
  \item \textbf{Forward pass}: $\hat{u} = \mathcal{N}(t,x;\bm\theta)$.
  \item \textbf{First derivatives}: Compute $\nabla_{(t,x)} \hat{u}$
        via a single \texttt{autograd.grad} call.  Since $\hat{u}$
        is a vector (one value per sample), we pass
        \texttt{grad\_outputs = \bm{1}}$ (the vector-Jacobian product
        $\bm{1}^\top \frac{\partial \hat{u}}{\partial (t,x)}$):
        \[
          \frac{\partial \hat{u}}{\partial t} = \text{grads}[:,0],
          \qquad
          \frac{\partial \hat{u}}{\partial x} = \text{grads}[:,1]
        \]
  \item \textbf{Second derivative}: Differentiate $u_x$ again
        w.r.t.\ the input tensor to get $u_{xx}$:
        \[
          \frac{\partial^2 \hat{u}}{\partial x^2}
          = \frac{\partial}{\partial x}
            \!\left(\frac{\partial \hat{u}}{\partial x}\right)
        \]
  \item \textbf{Assemble}: $f = u_t + \hat{u}\,u_x - \nu\,u_{xx}$.
\end{enumerate}

The flag \texttt{create\_graph=True} is essential: it tells PyTorch to
record the differentiation operation itself in the computational graph, so
that we can later backpropagate through the residual to update $\bm\theta$.

\subsection*{Code}

\begin{lstlisting}
def output(N, X):
    Nout = N(X)
    u = Nout[:, 0:1]
    return u

def get_results(N, X):
    X.requires_grad_(True)
    u = output(N, X)

    # du/dt and du/dx
    grads = torch.autograd.grad(
        u, X,
        grad_outputs=torch.ones_like(u),
        create_graph=True,
        retain_graph=True)[0]
    u_t = grads[:, 0]
    u_x = grads[:, 1]

    # d^2u/dx^2
    u_xx = torch.autograd.grad(
        u_x, X,
        grad_outputs=torch.ones_like(u_x),
        create_graph=True,
        retain_graph=True)[0][:, 1]

    # PDE residual
    fu = u_t + u[:, 0] * u_x - nu * u_xx
    return u, fu
\end{lstlisting}


% ======================================================================
\section{Section 6: Loss Function}
\label{sec:loss}
% ======================================================================

\subsection*{Mathematical description}

The total PINN loss is a weighted sum of three mean-squared-error terms:

\begin{equation}
\boxed{
  \mathcal{L}(\bm\theta)
  = \underbrace{
      \frac{1}{N_{\mathrm{int}}} \sum_{i=1}^{N_{\mathrm{int}}}
      \bigl[f(t_i, x_i;\bm\theta)\bigr]^2
    }_{\mathcal{L}_{\mathrm{PDE}}}
  + \lambda_0
    \underbrace{
      \frac{1}{N_0} \sum_{j=1}^{N_0}
      \bigl[\hat{u}(0, x_j) - u_0(x_j)\bigr]^2
    }_{\mathcal{L}_{\mathrm{IC}}}
  + \lambda_B
    \underbrace{
      \Bigl(
        \frac{1}{N_b^L}\!\sum_k\!\bigl[\hat{u}(t_k, -1)\bigr]^2
        + \frac{1}{N_b^R}\!\sum_k\!\bigl[\hat{u}(t_k, +1)\bigr]^2
      \Bigr)
    }_{\mathcal{L}_{\mathrm{BC}}}
}
\label{eq:loss}
\end{equation}

where:
\begin{itemize}[nosep]
  \item $\mathcal{L}_{\mathrm{PDE}}$: enforces the Burgers' equation in
        the domain interior.
  \item $\mathcal{L}_{\mathrm{IC}}$: enforces the initial condition
        $u(0,x) = -\sin(\pi x)$.
  \item $\mathcal{L}_{\mathrm{BC}}$: enforces the Dirichlet boundary
        conditions $u(t,\pm 1) = 0$.
  \item $\lambda_0 = 5$ and $\lambda_B = 5$ weight the data terms
        relative to the PDE residual.  Higher weights force stricter
        enforcement of boundary/initial data.
\end{itemize}

\subsection*{Code}

\begin{lstlisting}
loss_function = nn.MSELoss()
lam0 = 5.0   # IC weight
lamB = 5.0   # BC weight

def loss(fu, u0, u0pinn, ul, ulpinn, ur, urpinn):
    Ntot = fu.shape[0]
    zeros = torch.zeros((Ntot, 1),
        dtype=torch.get_default_dtype(),
        device=fu.device)
    fu_col = fu.reshape(-1, 1)

    loss_value = (
        loss_function(fu_col, zeros)
        + lam0 * loss_function(u0, u0pinn)
        + lamB * (loss_function(ul, ulpinn)
                  + loss_function(ur, urpinn))
    )
    return loss_value
\end{lstlisting}


% ======================================================================
\section{Section 7: Gradient Computation and Adam Training Step}
\label{sec:adam_step}
% ======================================================================

\subsection*{Mathematical description}

During Phase~1 (Adam), at each iteration we:

\begin{enumerate}[nosep]
  \item Evaluate the loss $\mathcal{L}(\bm\theta)$ from Eq.~\eqref{eq:loss}.
  \item Compute the gradient $\nabla_{\bm\theta} \mathcal{L}$ via
        backpropagation (\texttt{loss.backward()}).
  \item Update parameters using Adam:
\end{enumerate}

\begin{align}
  \bm{m}_t &= \beta_1\,\bm{m}_{t-1} + (1-\beta_1)\,\bm{g}_t
    & &\text{(first moment estimate)} \\
  \bm{v}_t &= \beta_2\,\bm{v}_{t-1} + (1-\beta_2)\,\bm{g}_t^2
    & &\text{(second moment estimate)} \\
  \hat{\bm{m}}_t &= \bm{m}_t / (1-\beta_1^t)
    & &\text{(bias correction)} \\
  \hat{\bm{v}}_t &= \bm{v}_t / (1-\beta_2^t) \\
  \bm\theta_{t+1} &= \bm\theta_t
    - \eta\,\hat{\bm{m}}_t / (\sqrt{\hat{\bm{v}}_t} + \epsilon)
    & &\text{(parameter update)}
\label{eq:adam}
\end{align}
with $\eta = 5\times 10^{-3}$, $\beta_1 = 0.99$, $\beta_2 = 0.999$,
$\epsilon = 10^{-20}$.

\subsection*{Code}

\begin{lstlisting}
def grads(N, X, X0, Xb):
    for p in N.parameters():
        p.grad = None

    X  = X.clone().detach().requires_grad_(True)
    X0 = X0.clone().detach().requires_grad_(True)

    _, fu   = get_results(N, X)
    u0      = uinit(X0)
    u0pinn  = output(N, X0)
    ul      = uleft(Xb)
    ur      = uright(Xb)
    mask_l  = (Xb[:, 1] == x0)
    mask_r  = (Xb[:, 1] == xf)
    ulpinn  = output(N, Xb[mask_l])
    urpinn  = output(N, Xb[mask_r])

    loss_value = loss(fu, u0, u0pinn,
                      ul, ulpinn, ur, urpinn)
    loss_value.backward()

    gradsN = [p.grad for p in N.parameters()]
    return gradsN, loss_value


def training(N, X, X0, Xb, optimizer):
    _, loss_value = grads(N, X, X0, Xb)
    optimizer.step()
    return loss_value
\end{lstlisting}


% ======================================================================
\section{Section 8: Generate Initial Training Data}
\label{sec:data}
% ======================================================================

\subsection*{Mathematical description}

Before training begins, we sample the three point sets:
\begin{itemize}[nosep]
  \item $N_{\mathrm{int}} = 8000$ interior points
        $(t_i, x_i) \sim \mathcal{U}([0,1]\times[-1,1])$
  \item $N_0 = 500$ initial-condition points $(0, x_j)$
  \item $N_b = 500$ boundary points $(t_k, \pm 1)$
\end{itemize}

These are re-sampled periodically during training using RAD (Section~\ref{sec:sampling}).

\subsection*{Code}

\begin{lstlisting}
X  = generate_inputs(Nint)   # 8000 interior pts
X0 = initial_points(N0)      # 500 IC pts
Xb = boundary_points(Nb)     # 500 BC pts
\end{lstlisting}


% ======================================================================
\section{Section 9: Load Reference Solution}
\label{sec:reference}
% ======================================================================

\subsection*{Mathematical description}

A high-fidelity reference solution $u_{\mathrm{ref}}(t,x)$ is loaded from
a MATLAB \texttt{.mat} file.  This is the ``ground truth'' obtained by a
conventional numerical method (e.g., spectral or finite-difference solver).

The relative $\ell^2$ error is computed as:
\begin{equation}
  e_{\mathrm{rel}} = \frac{\|u_{\mathrm{ref}} - \hat{u}\|_2}{\|\hat{u}\|_2 + \epsilon}
\label{eq:rel_error}
\end{equation}
with $\epsilon = 10^{-12}$ for numerical stability.

\subsection*{Code}

\begin{lstlisting}
mat    = scipy.io.loadmat("burgers_canonical.mat")
u_ref  = mat["usol"]
t_star = mat["t"].ravel()
x_star = mat["x"].ravel()
x_star, t_star = np.meshgrid(x_star, t_star)

X_star = torch.as_tensor(
    np.hstack((t_star.reshape(-1,1),
               x_star.reshape(-1,1))),
    dtype=torch.get_default_dtype(),
    device=device)
\end{lstlisting}


% ======================================================================
\section{Section 10: Phase 1 --- Adam Training}
\label{sec:adam}
% ======================================================================

\subsection*{Mathematical description}

Phase~1 runs $1{,}000$ Adam iterations. The learning rate follows an
exponential decay schedule:
\begin{equation}
  \eta_t = \eta_0 \cdot 0.98^{t/1000}
  = 5\times10^{-3} \cdot 0.98^{t/1000}
\label{eq:lr_schedule}
\end{equation}

Every $N_{\mathrm{change}} = 500$ epochs, the collocation points are
resampled using RAD (Eq.~\eqref{eq:rad}), which concentrates points
in regions of high PDE residual.  This adaptive refinement is critical
for problems with steep gradients or shocks.

Every $N_{\mathrm{print}} = 100$ epochs, the relative $\ell^2$ error
(Eq.~\eqref{eq:rel_error}) is computed against the reference solution
for monitoring.

\emph{Purpose:} Adam is fast per-iteration but converges slowly to
high accuracy.  It gets the network into a good basin of attraction
from which the more powerful second-order optimizer can take over.

\subsection*{Code}

\begin{lstlisting}
rel_adam = []
optimizer = torch.optim.Adam(
    N.parameters(), lr=5e-3,
    betas=(0.99, 0.999), eps=1e-20)
adam_scheduler = torch.optim.lr_scheduler.LambdaLR(
    optimizer,
    lr_lambda=lambda step: 0.98 ** (step / 1000.0))

adam_losses = []
adam_t0 = perf_counter()

for epoch in range(Nepochs_ADAM):
    if (epoch + 1) % Nchange == 0:
        X  = adaptive_rad(N, Nint, rad_args)
        X0 = initial_points(N0)
        Xb = boundary_points(Nb)

    loss_value = training(N, X, X0, Xb, optimizer)
    adam_losses.append(
        float(loss_value.detach().cpu().item()))

    if (epoch + 1) % Nprint == 0:
        print(template.format(
            epoch + 1, adam_losses[-1]))
        N.eval()
        with torch.no_grad():
            u_pred = output(N, X_star) \
                .detach().cpu().numpy() \
                .reshape(u_ref.shape)
        N.train()
        num = np.linalg.norm(u_ref - u_pred)
        den = np.linalg.norm(u_pred) + 1e-12
        rel = num / den
        rel_adam.append((epoch + 1, float(rel)))

adam_time_sec = perf_counter() - adam_t0
\end{lstlisting}


% ======================================================================
\section{Section 11: Nested Tensor Utility}
\label{sec:nested}
% ======================================================================

\subsection*{Mathematical description}

SciPy's \texttt{minimize} works with a single flat vector
$\bm\theta \in \mathbb{R}^{1341}$, but PyTorch stores parameters as
structured tensors (weight matrices and bias vectors).  The
\texttt{nested\_tensor} utility converts between the two representations.

For \texttt{layer\_dims = (2, 20, 20, 20, 20, 1)}, the flat vector is
unpacked into 10 arrays:
\[
  \bm\theta = [\,
    \underbrace{\bm{W}^{[1]}}_{2\times 20},\;
    \underbrace{\bm{b}^{[1]}}_{20},\;
    \underbrace{\bm{W}^{[2]}}_{20\times 20},\;
    \underbrace{\bm{b}^{[2]}}_{20},\;
    \ldots,\;
    \underbrace{\bm{W}^{[5]}}_{20\times 1},\;
    \underbrace{\bm{b}^{[5]}}_{1}
  \,]
\]

The flat vector has length:
\[
  |\bm\theta| = \sum_{\ell=1}^{5}
  \bigl(d_{\ell-1} \cdot d_\ell + d_\ell\bigr)
  = 40{+}20 + 400{+}20 + 400{+}20 + 400{+}20 + 20{+}1 = 1341
\]

\subsection*{Code}

\begin{lstlisting}
def nested_tensor(grad, layer_dims,
                  train_activations=False,
                  bias=True):
    grad = np.asarray(grad).ravel()
    if not train_activations:
        if bias:
            temp = [None]*(2*len(layer_dims) - 2)
        else:
            temp = [None]*(2*len(layer_dims) - 3)
        index = 0
        for i in range(len(temp)):
            if i % 2 == 0:
                k = i // 2
                fan_in  = layer_dims[k]
                fan_out = layer_dims[k + 1]
                size = fan_in * fan_out
                temp[i] = grad[index:index+size] \
                    .reshape(fan_in, fan_out)
                index += size
            else:
                k = i - (i // 2)
                bsz = layer_dims[k]
                temp[i] = grad[index:index+bsz]
                index += bsz
        return temp
\end{lstlisting}


% ======================================================================
\section{Section 12: SciPy-Compatible Loss + Gradient Wrapper}
\label{sec:wrapper}
% ======================================================================

\subsection*{Mathematical description}

SciPy's \texttt{minimize(fun, x0, jac=True)} expects a callable:
\[
  \texttt{fun}(\bm\theta) \;\to\;
  \bigl(\mathcal{L}(\bm\theta),\;
        \nabla_{\bm\theta}\mathcal{L}\bigr)
  \qquad
  \bm\theta \in \mathbb{R}^{1341},\;
  \nabla_{\bm\theta}\mathcal{L} \in \mathbb{R}^{1341}
\]

The wrapper performs these steps:
\begin{enumerate}[nosep]
  \item Receive $\bm\theta$ as a flat \texttt{numpy} array from SciPy.
  \item Convert to a PyTorch tensor and load into the network via
        \texttt{vector\_to\_parameters}.
  \item Run the PyTorch forward/backward pass to get $\mathcal{L}$ and
        $\nabla_{\bm\theta}\mathcal{L}$.
  \item Flatten the gradient list and return both as \texttt{numpy} arrays.
\end{enumerate}

An optional \textbf{power transform} is supported:
$\mathcal{L}_{\mathrm{eff}} = \mathcal{L}^{1/p}$.
When $p=1$ (default), the loss is used directly.  When $p=2$, the
optimizer minimizes $\sqrt{\mathcal{L}}$, which can improve conditioning.

\subsection*{Code}

\begin{lstlisting}
power = 1.0

def loss_and_gradient_torch(N, X, X0, Xb,
                            power=power):
    X = X.clone().detach().requires_grad_(True)
    _, fu = get_results(N, X)
    u0     = uinit(X0)
    u0pinn = output(N, X0)
    ul     = uleft(Xb)
    mask_l = (Xb[:, 1] == x0)
    ulpinn = output(N, Xb[mask_l])
    ur     = uright(Xb)
    mask_r = (Xb[:, 1] == xf)
    urpinn = output(N, Xb[mask_r])

    loss_value = loss(fu, u0, u0pinn,
                      ul, ulpinn, ur, urpinn)
    loss_root = (loss_value
        if power == 1.0
        else loss_value ** (1.0 / power))

    params = list(N.parameters())
    gradsN = torch.autograd.grad(
        loss_root, params,
        create_graph=False,
        retain_graph=False,
        allow_unused=False)
    return loss_root, gradsN


def loss_and_gradient(weights, N, X, X0, Xb,
                      layer_dims=None):
    first_param = next(N.parameters())
    device = first_param.device
    dtype  = first_param.dtype
    w_tensor = torch.as_tensor(
        weights, dtype=dtype, device=device)
    with torch.no_grad():
        vector_to_parameters(
            w_tensor, N.parameters())
    loss_val, grads_list = \
        loss_and_gradient_torch(
            N, X, X0, Xb, power=power)
    grads_flat = torch.cat(
        [g.reshape(-1) for g in grads_list])
    return (float(loss_val.detach().cpu().item()),
            grads_flat.detach().cpu().numpy())
\end{lstlisting}


% ======================================================================
\section{Section 13: Phase 2 Setup --- Self-Scaled BFGS}
\label{sec:bfgs_setup}
% ======================================================================

\subsection*{Mathematical description}

\subsubsection*{Why switch from Adam to BFGS?}

Adam is a \emph{first-order} method: each step uses only $\nabla\mathcal{L}$.
Near the solution, where the loss landscape is nearly quadratic, the
convergence rate of first-order methods is only \emph{linear} (geometric
with a rate $\approx 1 - 2/\kappa$, where $\kappa$ is the condition number
of the Hessian).

In contrast, \textbf{BFGS} is a \emph{quasi-Newton} method that
builds an approximate inverse Hessian $H_k \approx
[\nabla^2 \mathcal{L}]^{-1}$, achieving \emph{superlinear} convergence
near the solution.

\subsubsection*{Standard BFGS update}

At iteration $k$, after a line search determines step size $\alpha_k$:
\begin{align}
  \bm{s}_k &= \alpha_k\,\bm{p}_k = \bm\theta_{k+1} - \bm\theta_k \\
  \bm{y}_k &= \nabla\mathcal{L}_{k+1} - \nabla\mathcal{L}_k \\
  \rho_k   &= \frac{1}{\bm{y}_k^\top \bm{s}_k}
\end{align}
The inverse Hessian is updated:
\begin{equation}
  H_{k+1} = H_k
  - \rho_k\,(H_k \bm{y}_k\,\bm{s}_k^\top + \bm{s}_k\,\bm{y}_k^\top H_k)
  + \rho_k\,(1 + \bm{y}_k^\top H_k\,\bm{y}_k\,\rho_k)\,
    \bm{s}_k\,\bm{s}_k^\top
\label{eq:bfgs}
\end{equation}

\subsubsection*{Self-Scaled Broyden (SSBroyden2)}

The \textbf{self-scaled} variants (from the paper) introduce a scaling
factor $\tau_k$ that multiplies $H_k$ \emph{before} the rank-2 update,
plus a parameter $\theta_k$ that interpolates between DFP and BFGS
within the Broyden family:

\begin{equation}
  H_{k+1} = \frac{1}{\tau_k}\left(
    H_k
    - \frac{H_k\bm{y}_k\bm{y}_k^\top H_k}{\bm{y}_k^\top H_k\bm{y}_k}
    + \phi_k\,(\bm{y}_k^\top H_k\bm{y}_k)\,\bm{v}_k\bm{v}_k^\top
  \right)
  + \rho_k\,\bm{s}_k\bm{s}_k^\top
\label{eq:ssbroyden}
\end{equation}
where:
\begin{align}
  h_k &= \rho_k\,\bm{y}_k^\top H_k\bm{y}_k,
  \qquad
  b_k = -\alpha_k\,\rho_k\,\bm{s}_k^\top(\nabla\mathcal{L}_{k+1} - \bm{y}_k) \\
  a_k &= b_k\,h_k - 1 \label{eq:ak}\\
  \bm{v}_k &= \rho_k\,\bm{s}_k
    - \frac{H_k\,\bm{y}_k}{\bm{y}_k^\top H_k\,\bm{y}_k}
  \label{eq:vk}\\
  \phi_k &= \frac{1 - \theta_k}{1 + a_k\,\theta_k}
  \label{eq:phik}
\end{align}

In \texttt{SSBroyden2}, $\theta_k$ is chosen with a clamped formula:
\begin{equation}
  \theta_k = \max\!\left(\theta_k^{-},\;
    \min\!\left(\theta_k^{+},\; \frac{1-b_k}{b_k}\right)\right)
\label{eq:theta_clamp}
\end{equation}
where $\theta_k^{\pm}$ are bounds derived from positive-definiteness
constraints.  The scaling factor $\tau_k$ is computed as:
\begin{equation}
  \sigma_k = 1 + \theta_k\,a_k,
  \qquad
  \sigma_k^{1/(1-N)} = |\sigma_k|^{1/(1-N)}
\end{equation}
\begin{equation}
  \tau_k = \begin{cases}
    \min\!\bigl(\rho_k^{(k)}\,\sigma_k^{1/(1-N)},\; \sigma_k\bigr)
      & \text{if } \theta_k \le 0 \\
    \rho_k^{(k)}\,\min\!\bigl(\sigma_k^{1/(1-N)},\; 1/\theta_k\bigr)
      & \text{if } \theta_k > 0
  \end{cases}
\label{eq:tauk}
\end{equation}
with $\rho_k^{(k)} = \min(1, 1/b_k)$.

This scaling keeps the condition number of $H_k$ controlled and
improves convergence over standard BFGS.

\subsection*{Setup code}

\begin{lstlisting}
Nbfgs   = 10000
Nchange = 500
Nprint  = 100
n_ckpts = Nbfgs // Nprint

lossbfgs        = np.zeros(n_ckpts)
validation_list = np.zeros(n_ckpts)
error_list      = np.zeros(n_ckpts)

Nt = 300; Nx = 300
Xtest, t, x = generate_validation(Nt, Nx)

initial_weights = parameters_to_vector(
    [p.detach() for p in N.parameters()]
).cpu().numpy()

cont = 0   # global BFGS iteration counter
\end{lstlisting}


% ======================================================================
\section{Section 14: BFGS Optimizer Configuration}
\label{sec:bfgs_config}
% ======================================================================

\subsection*{Mathematical description}

The optimizer is configured as:
\begin{itemize}[nosep]
  \item \texttt{method = "BFGS"}: selects the BFGS family in SciPy.
  \item \texttt{method\_bfgs = "SSBroyden2"}: selects the self-scaled
        Broyden variant~2 from the patched \texttt{\_optimize.py}.
  \item \texttt{initial\_scale = False}: do \emph{not} apply the
        Nocedal--Wright initial scaling $H_0 \leftarrow H_0 /
        (\rho_0\,\|\bm{y}_0\|^2)$ at the first iteration.
  \item The initial inverse Hessian is $H_0 = I_{1341\times 1341}$,
        meaning the first BFGS step is simply a gradient descent step:
        $\Delta\bm\theta = -I\cdot\nabla\mathcal{L} = -\nabla\mathcal{L}$.
\end{itemize}

A \textbf{callback} function is registered that, every 100 iterations:
\begin{enumerate}[nosep]
  \item Records the training loss.
  \item Evaluates the network on a $300\times 300$ validation grid.
  \item Computes the relative $\ell^2$ error vs.\ the reference solution.
\end{enumerate}

\subsection*{Code}

\begin{lstlisting}
method      = "BFGS"
method_bfgs = "SSBroyden2"
initial_scale = False

H0 = np.eye(initial_weights.size,
             dtype=np.float64)

def callback(*, intermediate_result):
    global cont, lossbfgs, error_list
    cont += 1
    if (cont % Nprint) != 0:
        return
    idx = cont // Nprint - 1
    loss_value = float(
        intermediate_result.fun ** power)
    lossbfgs[idx] = loss_value

    # Validation error
    _, futest = get_results(N, Xtest.to(device))
    with torch.no_grad():
        u_pred = output(N, X_star) \
            .reshape(x_star.shape)
        u_ref_t = torch.as_tensor(u_ref)
        v1 = loss_function(u_pred, u_ref_t)
        v2 = loss_function(
            futest.reshape(-1,1),
            torch.zeros_like(
                futest.reshape(-1,1)))
        validation_list[idx] = (v1+v2).item()
    u_np = u_pred.detach().cpu().numpy()
    error_list[idx] = (
        np.linalg.norm(u_ref - u_np)
        / (np.linalg.norm(u_np) + 1e-12))
\end{lstlisting}


% ======================================================================
\section{Section 15: Phase 2 --- Self-Scaled Broyden Training Loop}
\label{sec:bfgs_loop}
% ======================================================================

\subsection*{Mathematical description}

The second-order optimization is performed in a \textbf{batched loop}:

\begin{enumerate}
  \item Call \texttt{scipy.optimize.minimize(\ldots, maxiter=500)}.
        This runs up to 500 BFGS iterations on the \emph{current}
        collocation points.
  \item \textbf{Extract and warm-start $H_k$}:
        \begin{itemize}[nosep]
          \item Retrieve $H_k = \texttt{result.hess\_inv}$.
          \item Symmetrize: $H_k \leftarrow \frac{1}{2}(H_k + H_k^\top)$.
          \item Check positive-definiteness via Cholesky decomposition.
                If $H_k$ has lost positive-definiteness (numerical
                accumulation of errors), reset $H_k \leftarrow I$.
        \end{itemize}
  \item \textbf{Resample collocation points} using RAD
        (Eq.~\eqref{eq:rad}).  This changes the loss landscape, which is
        why we run in batches rather than a single long call.
  \item Repeat until the total iteration budget (10\,000) is exhausted.
\end{enumerate}

\subsubsection*{Why batch rather than one long run?}

The loss function $\mathcal{L}(\bm\theta)$ depends on the particular
choice of collocation points.  If we never resample, the network may
overfit to those specific points  and fail to generalize.  Periodic
resampling (especially with RAD) acts as a form of stochastic
regularization and ensures the network learns the PDE everywhere in
the domain, not just at fixed locations.

The inverse Hessian $H_k$ is \emph{carried forward} (warm-started)
between batches, preserving the curvature information accumulated so
far.  However, when the collocation points change, the loss surface
changes as well, so the Hessian approximation becomes slightly stale.
The self-scaling factor $\tau_k$ in the Broyden update
(Eq.\eqref{eq:ssbroyden}) helps correct for this by adaptively
rescaling $H_k$ at each step.

\subsection*{Code}

\begin{lstlisting}
bfgs_t0 = perf_counter()

while cont < Nbfgs:
    result = minimize(
        loss_and_gradient,
        initial_weights,
        args=(N, X, X0, Xb, layer_dims),
        method=method,
        jac=True,
        options={
            'maxiter': Nchange,
            'gtol': 0,
            'hess_inv0': H0,
            'method_bfgs': method_bfgs,
            'initial_scale': initial_scale},
        tol=0,
        callback=callback)

    # Warm-start: carry forward weights + Hessian
    initial_weights = result.x
    H0 = result.hess_inv
    H0 = 0.5 * (H0 + H0.T)   # symmetrize

    # Check positive-definiteness
    try:
        cholesky(H0)
    except LinAlgError:
        H0 = np.eye(len(initial_weights),
                     dtype=np.float64)

    # Resample with RAD
    X  = adaptive_rad(N, Nint, rad_args)
    X0 = initial_points(N0)
    Xb = boundary_points(Nb)

    # Monitor loss on new points
    _, fu  = get_results(N, X)
    u0     = uinit(X0)
    u0pinn = output(N, X0)
    ul     = uleft(Xb)
    ulpinn = output(N, Xb[Xb[:,1]==x0])
    ur     = uright(Xb)
    urpinn = output(N, Xb[Xb[:,1]==xf])
    loss_value = loss(fu, u0, u0pinn,
                      ul, ulpinn, ur, urpinn)
    initial_scale = False

bfgs_time_sec = perf_counter() - bfgs_t0
\end{lstlisting}


% ======================================================================
\section{Summary of the Two-Phase Training Pipeline}
% ======================================================================

\begin{center}
\renewcommand{\arraystretch}{1.4}
\begin{tabular}{@{}lcc@{}}
\toprule
& \textbf{Phase 1: Adam} & \textbf{Phase 2: SSBroyden2} \\
\midrule
Type       & First-order     & Quasi-Newton (second-order) \\
Iterations & 1\,000          & 10\,000 \\
Per-step cost & $O(n)$       & $O(n^2)$ (Hessian-vector product) \\
Convergence & Linear         & Superlinear \\
Key idea   & Cheap gradient steps & Curvature-informed steps \\
Scaling    & ---             & Self-scaled $\tau_k$ controls $\kappa(H_k)$ \\
Resampling & RAD every 500   & RAD every 500 \\
\bottomrule
\end{tabular}
\end{center}

The overall flow is:
\[
  \underbrace{\text{Adam (1000 iters)}}_{\text{get into a good basin}}
  \;\longrightarrow\;
  \underbrace{\text{SSBroyden2 (10000 iters)}}_{\text{converge to high accuracy}}
\]

The self-scaled Broyden update (Eq.~\eqref{eq:ssbroyden}) with the
$\tau_k$ scaling (Eq.~\eqref{eq:tauk}) is what distinguishes this
implementation from standard BFGS and is the core contribution of the
paper.

% ======================================================================
\appendix
\section{Notation Reference}
% ======================================================================

\begin{center}
\renewcommand{\arraystretch}{1.3}
\begin{tabular}{@{}cl@{}}
\toprule
\textbf{Symbol} & \textbf{Meaning} \\
\midrule
$u(t,x)$          & Solution of Burgers' equation \\
$\hat{u}(t,x;\bm\theta)$ & Neural network approximation \\
$\bm\theta$        & All trainable parameters (flat vector, $\in\mathbb{R}^{1341}$) \\
$\nu$              & Viscosity ($0.01/\pi$) \\
$f(t,x;\bm\theta)$ & PDE residual: $u_t + u\,u_x - \nu\,u_{xx}$ \\
$\mathcal{L}$      & Total loss (PDE + IC + BC) \\
$\lambda_0, \lambda_B$ & Loss weights for IC and BC (both $= 5$) \\
$H_k$              & Inverse Hessian approximation at iteration $k$ \\
$\tau_k$           & Self-scaling factor \\
$\theta_k$         & Broyden family interpolation parameter \\
$\bm{s}_k, \bm{y}_k$ & Step and gradient difference vectors \\
$\rho_k$           & $1 / (\bm{y}_k^\top \bm{s}_k)$ \\
$e_{\mathrm{rel}}$ & Relative $\ell^2$ error \\
\bottomrule
\end{tabular}
\end{center}

\end{document}
